\documentclass[11pt]{article}
\usepackage[margin=25mm]{geometry}
\usepackage{fontspec}
\setmainfont{Latin Modern Roman}
\usepackage{amsmath,amssymb,mathtools}
\usepackage{unicode-math}
\setmathfont{Latin Modern Math}
\usepackage{xurl,hyperref}
\hypersetup{colorlinks=true,urlcolor=blue}
\setcounter{secnumdepth}{0}
\setlength{\parindent}{0pt}
\setlength{\parskip}{5pt}
\setlength{\emergencystretch}{3em}
\newcommand{\sref}[2]{\hyperlink{source:#1:#2}{[S#2]}}
\newcommand{\cataloguescope}[1]{\paragraph{Recorded target.}#1}
\newcommand{\R}{\mathbb R}
\newcommand{\C}{\mathbb C}
\newcommand{\N}{\mathbb N}
\newcommand{\Z}{\mathbb Z}
\newcommand{\Q}{\mathbb Q}
\newcommand{\F}{\mathbb F}
\newcommand{\eps}{\varepsilon}
\begin{document}
% UnsolvedMath ID 30003147; source catalogue code OWR-14608-002
\setcounter{section}{30003146}
\section{Quantized Transport in Interacting Disordered Systems}\label{30003147}
{\small UnsolvedMath ID 30003147 · Mathematical Physics}
\subsection{Definitions and mathematical statement}

\paragraph{Shared notation.}$\mathbb N=\{1,2,\ldots\}$; the standard number fields have their usual meanings. All additional conventions are stated locally.


\paragraph{Definitions and conventions.}
The Hall conductivity is the transverse charge response to an applied electric field. A Hall plateau is a parameter interval on which this response takes a fixed quantized value. Disorder removes translation invariance and interactions couple particle degrees of freedom. Universal critical conductivity concerns response at transitions between such phases, requiring an identified class of microscopic Hamiltonians and a scaling limit.

\paragraph{Target retained from the canonical source catalogue.}
\paragraph{Catalogue wording retained for traceability.}
Develop a theory of quantized charge transport, including sharply quantized Hall plateaux and universal critical conductivities, for interacting disordered systems.

\paragraph{Source scope and qualification.}
This is a theory-development programme, not a universal quantization theorem for arbitrary interacting disordered Hamiltonians. The symmetry, dimensional, temperature and localization regimes are not fixed in the extracted question.
\subsection{Short English statement}
Build a theory of Hall plateaux and critical charge transport that includes both interactions and disorder.
\subsection{Research attempt}
\paragraph{Model, date and investigated scope.}
OpenAI GPT-6 Astra (configured model \texttt{gpt-6-astra}), ultra reasoning; generated 11 October 2026. Proves integer curvature flux and its stability in a conditional finite-dimensional gapped twist family; does not prove Hall plateaux or universal critical conductivity in interacting disordered matter.

\paragraph{Sources and literature check.}
The complete supplied definition, including source qualifications and the dated literature assessment, was reviewed. The calculation below is an elementary restricted attempt motivated by [S2]; it does not claim novelty or a fresh literature-status determination.

\paragraph{Context and investigated scope.} The complete supplied definition, including source qualifications and the dated literature assessment, was reviewed. The calculation below is an elementary restricted attempt motivated by [S2]; it does not claim novelty or a fresh literature-status determination.

\paragraph{A finite-dimensional quantization mechanism.}
Consider a smooth family of finite-dimensional Hermitian matrices $H(\theta_1,\theta_2)$, periodic with period $2\pi$ in each twist variable, with a simple ground eigenvalue separated from the rest of the spectrum everywhere. To make the argument explicit, assume a normalized ground eigenvector $u$ has been chosen smoothly on the cylinder $0\leq\theta_2\leq2\pi$, periodically in $\theta_1$. At the identified boundary circles,
\[
 u(\theta_1,2\pi)=e^{i\chi(\theta_1)}u(\theta_1,0),
 \qquad \chi(2\pi)-\chi(0)=2\pi n,\quad n\in\mathbb Z.
\]
The phase has a continuous real lift along $0\leq\theta_1\leq2\pi$; the endpoint difference is an integer multiple of $2\pi$ because its exponential is periodic.

Set $A=i\langle u,du\rangle$ and $F=dA$. Normalization makes $A$ real. The seam condition gives $A_{\rm top}=A_{\rm bottom}-d\chi$. Stokes' formula, oriented by $d\theta_1\wedge d\theta_2$, yields
\[
 C:=\frac1{2\pi}\int_{[0,2\pi]^2}F
 =\frac1{2\pi}\int_0^{2\pi}
       (A_{1,\rm bottom}-A_{1,\rm top})\,d\theta_1
 =n.
\]
The vertical boundary contributions cancel by periodicity. A smooth change of local phase changes $A$ by an exact form and leaves $F$ fixed. This is the elementary line-bundle quantization underlying the candidate transport invariant.

\paragraph{What produces a plateau in this model?}
If the matrices depend smoothly on an additional parameter $s$ in an interval, the gap remains positive, and the ground projections stay smooth, the curvature integral is continuous in $s$. Since it is integer valued, it is constant. Thus a gap-preserving path cannot change this finite-system invariant. The argument permits disorder in the matrices; it uses no translation symmetry.

\paragraph{The missing transport statement.}
An integer curvature integral is not by itself the physical Hall response at a fixed boundary twist. The identification with an appropriate Kubo response, removal of twist averaging, the thermodynamic and low-frequency limits, and plateau stability under a mobility gap require additional estimates. A uniform spectral gap is a strong hypothesis and is not established here for any source-specific interacting disordered ensemble. At a transition the gap hypothesis can fail, and integer continuity supplies no value or universality theorem for critical conductivity. The source's symmetry, dimension, temperature and localization qualifications, and its already established partial results, are therefore preserved.

\paragraph{Verification and adversarial audit.}
Checked connection gauge signs, orientation of the seam integral, integer winding, the gap hypothesis, and failure of this argument to determine a critical conductivity. Model self-audit only; no independent human proof checking. Prompts A and B were applied to this canonical record and its qualified source scope. Checked connection gauge signs, orientation of the seam integral, integer winding, the gap hypothesis, and failure of this argument to determine a critical conductivity. Retained unresolved status and zero original-target progress; Identify the finite-system curvature invariant with the physical infinite-volume transport coefficient under a specified localization regime and control the necessary limits.

Coordinator checked the seam phase, oriented Stokes formula and integrality of the curvature flux, and the conditional constancy along a gapped matrix path. This additional check was model-assisted, not independent human verification.

\paragraph{Tools, process and contributions.}
Python 3 was used for file preparation and structural checks; no numerical experiment or theorem prover was used. Contribution prompt A was followed by prompt B. GPT-6 Astra with ultra reasoning drafted the research and performed its stated self-audit; any additional coordinator contribution is identified in the verification and audit. A GPT-6-based Codex coordinator assembled the manuscript and checked repository integration; its exact variant and reasoning configuration were not exposed. No human mathematical contribution, independent human proof check or formal proof certification is supplied. Complete derivations and the consulted source references are included here; private reasoning transcripts are not published.

\paragraph{Final status and first remaining gap.}
\textbf{UNRESOLVED}. Identify the finite-system curvature invariant with the physical infinite-volume transport coefficient under a specified localization regime and control the necessary limits.

\section{Completion Estimate}
\noindent\textbf{COMPLETION: 0\%}
The calculation is an auxiliary or restricted attempt and establishes no newly verified portion of the full source target. Identify the finite-system curvature invariant with the physical infinite-volume transport coefficient under a specified localization regime and control the necessary limits.

\subsection{Sources}
\begin{itemize}
\item[\hypertarget{source:30003147:1}{S1}] ULAM.AI and the UnsolvedMath Contributors, UnsolvedMath, record 30003147 (OWR-14608-002). Catalogue curation under CC BY 4.0; definitions and English exposition added locally.
\url{https://huggingface.co/datasets/ulamai/UnsolvedMath}.
\item[\hypertarget{source:30003147:2}{S2}] The Renormalization Group (2016), Oberwolfach Reports, DOI 10.4171/owr/2016/26, pp. 1453--1515. Primary report accessed 2026-10-11; relevant PDF pages 45.
\url{https://doi.org/10.4171/owr/2016/26}.
\item[\hypertarget{source:30003147:3}{S3}] Mastropietro, Quantization of the interacting Hall conductivity (catalogue-listed context; not independently consulted in this preparation).
\url{https://pmc.ncbi.nlm.nih.gov/articles/PMC7405942/}.
\item[\hypertarget{source:30003147:4}{S4}] current primary work on interacting disordered edge transport (catalogue-listed context; not independently consulted in this preparation).
\url{https://link.springer.com/article/10.1007/s00220-026-05560-1}.
\end{itemize}
\label{end:30003147}
\end{document}
